\documentclass[12pt,a4paper]{article}
\usepackage[utf8]{inputenc}
\usepackage[T1]{fontenc}
\usepackage[bahasa]{babel}
\usepackage{mathptmx} % Times New Roman font
\usepackage{microtype} % Untuk optimasi typography dan spacing
\usepackage[none]{hyphenat} % Menonaktifkan hyphenation
\usepackage{graphicx}
\usepackage{amsmath}
\usepackage{amssymb}
\usepackage{setspace}
\usepackage{geometry}
\usepackage{tikz}
\usepackage{float}
\usepackage{enumitem}
\usepackage{hyperref}
\hypersetup{
    pdfencoding=auto,
    pdftitle={Bab 5: Hasil dan Pembahasan},
    pdfauthor={},
    pdfsubject={Tesis},
    colorlinks=true,
    linkcolor=black,
    citecolor=black,
    urlcolor=blue,
    pdfstartview=FitH,
    bookmarksopen=true,
    bookmarksnumbered=true
}
\usepackage{tocloft}
\usepackage{caption}
\usepackage{textcomp} % For special symbols
\usepackage{booktabs} % Untuk tabel yang lebih baik

% Aktifkan hyphenation Bahasa Indonesia dan optimasi spasi agar tidak keluar margin
\lefthyphenmin=2
\righthyphenmin=2
\pretolerance=1000
\tolerance=2000
\emergencystretch=3em
\hbadness=10000
\vbadness=10000
\hyphenpenalty=500
\exhyphenpenalty=500
\sloppy
% Aktifkan microtype untuk Bahasa Indonesia
\microtypesetup{activate=true,protrusion=true,expansion=true,tracking=true,kerning=true}
\microtypecontext{spacing=nonfrench}

% Pengaturan paragraf: tanpa indentasi dan jarak minimal antar paragraf
\setlength{\parindent}{0pt} % Menghilangkan indentasi awal paragraf
\setlength{\parskip}{0pt} % Tidak ada jarak antar paragraf

% Pengaturan jarak untuk subsection
\usepackage{titlesec}
\titleformat{\subsection}{\normalsize\bfseries}{\thesubsection}{1em}{}
\titlespacing*{\subsection}{0pt}{0pt}{0pt} % Tidak ada jarak sebelum atau setelah judul

% Pengaturan margin: kiri 4cm, lainnya 3cm
\geometry{
    a4paper,
    left=4cm,
    right=3cm,
    top=3cm,
    bottom=3cm
}

% Pengaturan spasi
\onehalfspacing

% Pengaturan heading
\renewcommand{\thesection}{\Roman{section}}
renewcommand{\thesubsection}{\thesection.\arabic{subsection}}
renewcommand{\thesubsubsection}{\thesubsection.\arabic{subsubsection}}

\begin{document}

% ============================================================
% BAB V: HASIL DAN PEMBAHASAN
% ============================================================

\vspace{2cm}
\begin{center}
{\fontsize{14}{16.8}\selectfont\textbf{BAB V. Hasil dan Pembahasan}}\\[1em]
\end{center}
\label{sec:hasil-pembahasan}
\addcontentsline{toc}{section}{BAB V HASIL DAN PEMBAHASAN}
\section{5}
\subsection{0}
\vspace{2em}

Bab ini menyajikan hasil komprehensif dari eksperimen restorasi dokumen berbasis Generative Adversarial Network dengan Diskriminator Dual-Modal yang telah dikembangkan. Pembahasan mencakup analisis mendalam terhadap hasil training dengan strategi Frozen Recognizer vs Joint Training, evaluasi performa menggunakan metrik objektif dan subjektif, perbandingan dengan metode state-of-the-art, serta validasi hipotesis penelitian melalui analisis statistik. Semua eksperimen dilakukan dengan kontrol yang ketat untuk memastikan validitas dan reproduktifitas hasil.

\end{document}